%% !TeX root=Proposal.tex
%% !TEX TS-program = XeLaTeX
%% Last version 1400.11.18
% نمونه پیشنهاد رساله (پروپوزال) دانشکده برق و کامپیوتر دانشگاه صنعتی اصفهان
%نظرات اصلاحی همه عزیزان را در ارتباط با قالب طراحی شده به دیده منت می‌نهیم و با کمال سپاس و قدردانی مورد استفاده قرار خواهیم داد (a.qazavi@ec.iut.ac.com).
%%%%%%%%%%%%%%%%%%%%%%%%%%%%%%%%%%%%%%%%%%%%%%%%%%%%%%%%%%%%%%%%%%%%%%%%%%%%%%%%%%%%%%%%%%%%%%%%%%%%%%%%%%%%%%%%%%%%%%%%%%%%
\documentclass[a4paper,fleqn]{article}
%در صورت نیاز به فراخوانی بسته‌های اضافی، آن‌ها را در فایل commands.tex و قبل از
%بسته زی‌پرشین فراخوانی کنید.
%% !TeX root=GUATThesis-main.tex
%% !TEX TS-program = XeLaTeX
%بسته های برای تایپ متون ریاضی و محیط های شماره دار
\usepackage{amsthm,amssymb,amsmath,amsfonts}			% Writing math
\usepackage{mathtools}
%% \usepackage{upgreek}  %%برای ایستاده نمایش دادن شکل ایستاده نمادهای یونانی
\usepackage[shortlabels]{enumitem}
%%%%%%%%%%%%%%%%%%%%%%%%%%%%%%%%%%%%%%%%%%%%%%%%%%%
% بسته‌ای برای تنطیم حاشیه‌های بالا، پایین، چپ و راست صفحه
\usepackage[top=40mm, bottom=20mm, left=20mm, right=20mm]{geometry}
%%%%%%%%%%%%%%%%%%%%%%%%%%%%%%%%%%%%%%%%%%%%%%%%%%%
% بسته ای برای رسم انواع مختلف نمودارها
\usepackage{tikz}
%%%%%%%%%%%%%%%%%%%%%%%%%%%%%%%%%%%%%%%%%%%%%%%%%%%
% بسته‌‌ای برای ظاهر شدن شکل‌ها و تصاویر متن
\usepackage{epsf,graphicx}								% Including graphics
\usepackage{caption, subcaption}
\usepackage{xcolor,colortbl}
\definecolor{green}{rgb}{0.5,1,0.5}
\usepackage{makecell}
%%%%%%%%%%%%%%%%%%%%%%%%%%%%%%%%%%%%%%%%%%%%%%%%%%
%دستوری برای وارد کردن واژه‌نامه انگلیسی به فارسی
\newcommand\persiangloss[2]{#1\dotfill\lr{#2}\\}
%\newcommand{\done}{\cellcolor{teal}5- اطلاعات رساله} %{0.9}
%\newcommand{\hcyan}[1]{{\color{teal} #1}}
%%%%%%%%%%%%%%%%%%%%%%%%%%%%%%%%%%%%%%%%%%%%%%%%%%%
\usepackage[stable,bottom]{footmisc}					% Move footnotes to the bottom of page
\usepackage{notoccite}% PREVENTS CITES IN CAPTIONS FROM MISNUMBERING YOUR REFERENCES 
\usepackage{cite}
\usepackage{url}
%%%%%%%%%%%%%%%%%%%%%%%%%%%%%%%%%%%%%%%%%%%%%%%%%%%
% برای مدیریت بر ساختار عناوین، هدر و ... مربوط به فصل، بخش و ...
\usepackage{emptypage}
\usepackage{titlesec}
\usepackage{setspace}									% Change line spacing
%%%%%%%%%%%%%%%%%%%%%%%%%%%%%%%%%%%%%%%%%%%%%%%%%%%
% مورد استفاده در تعاریف با استفاده از newcommand
%\usepackage{xspace}
%%%%%%%%%%%%%%%%%%%%%%%%%%%%%%%%%%%%%%%%%%%%%%%%%%%
%%%% بسته‌ و دستوراتی برای ایجاد لینک‌های رنگی با امکان جهش
\usepackage[colorlinks=true,linkcolor=Blue,urlcolor=Blue,citecolor=Blue]{hyperref}
% چنانچه قصد پرینت گرفتن نوشته خود را دارید، خط بالا را غیرفعال و  از دستور زیر استفاده کنید چون در صورت استفاده از دستور زیر‌‌، 
% لینک‌ها به رنگ سیاه ظاهر خواهند شد و برای پرینت گرفتن، مناسب‌تر خواهد بود.
%\usepackage[pagebackref=false]{hyperref}

%%%%%%%%%%%%%%%%%%%%%%%%%%%%%%%%%%%%%%%%%%%%%%%%%%%
%بسته های مربوط به ایجاد الگوریتم
\usepackage{algorithm,algcompatible}
%%%%%%%%%%%%%%%%%%%%%%%%%%%%%%%%%%%%%%%%%%%%%%%%%%%
%بسته هایی برای مدیریت کپشن در محیط های شناور مانند شکل و جدول و ...
\usepackage[textfont={small},labelfont={rm,small},format=hang,labelsep=quad,justification={centering},aboveskip=2pt,belowskip=4pt]{caption}
\usepackage{subcaption}
%%%%%%%%%%%%%%%%%%%%%%%%%%%%%%%%%%%%%%%%%%%%%%%%%%%

%%%%%%%%%%%%%%%%%%%%%%%%%%%%%%%%%%%%%%%%%%%%%%%%%

%%%%%%%%%%%%%%%%%%%%%%%%%%%%%%%%%%%%%%%%%%%%%%%%%%%
% بسته مورد نیاز برای نوشتن کدهای برنامه نویسی در نوشتار و رنگ
\usepackage{listings}
\usepackage{color}
\definecolor{Blue}{rgb}{0,0,.55}

%%%%%%%%%%%%%%%%%%%%%%%%%%%%%%%%%%%%%%%%%%%%%%%%%%%%%%%
\usepackage{multirow}
\usepackage{array}
\usepackage{pbox}
%تغییر اندازه سطرهای جداول
\renewcommand{\arraystretch}{1.2}
%%%%%%%%%%%%%%%%%%%%%%%%%%%%%%%%%%%%%%%%%%%%%%%%%%%%%%%%%%%%%%%%%%%%%%%%%%%%%%%%%%%%
\usepackage{roundbox,framed}
%%%%%%%%%%%%%%%%%%%%%%%%%%%%%%%%%%%%%%%%%%%%%%%%%%%%%%%%%%%%%%%%%%
%بسته‌ای برای ظاهر شدن «مراجع»  در فهرست مطالب
\usepackage{tocbibind}
\usepackage{zref-perpage}
\zmakeperpage[1]{footnote}
%%%%%%%%%%%%%%%%%%%%%%%%%%%%%%%%
\makeatletter
\renewcommand\tableofcontents{%
	\if@twocolumn
	\@restonecoltrue\onecolumn
	\else
	\@restonecolfalse
	\fi
	\centerline{\huge\bfseries\contentsname
		\@mkboth{%
			\MakeUppercase\contentsname}{\MakeUppercase\contentsname}}%
	\vskip 20\p@
	\@starttoc{toc}%
	\if@restonecol\twocolumn\fi
}
\makeatother
%%%%%%%%%%%%%%%%%%%%%%%%%%%%%%%%
%\usepackage{indentfirst}
% فراخوانی بسته زی‌پرشین و تعریف قلم فارسی و انگلیسی
%\usepackage[extrafootnotefeatures]{xepersian}
\usepackage{xepersian}
%دستور \usepackage{xepersian} باید بعد از همه بسته‌ها (پکیج‌ها)گذاشته شود یعنی زیپرشین آخرین بسته‌ای باید باشد که فراخوانی می‌کنید
%\settextfont[Scale=0.93]{HM XZar}
\settextfont{XB Zar}										% Persian font
\setlatintextfont{Times New Roman}
% تعریف قلم‌های فارسی و انگلیسی برای استفاده در بعضی از قسمت‌های متن
\defpersianfont\titr[Scale=1]{B Titr}
\defpersianfont\Naz[Scale=1]{B Nazanin}
\defpersianfont\traffic[Scale=1]{B Traffic}
% چنانچه می‌خواهید که اعداد داخل فرمول‌ها، فارسی باشد، دستور زیر را فعال کنید
%\setdigitfont{Yas}
%%%%%%%%%%%%%%%%%%%%%%%%%%%%%%%%%%%%%%%%%%

%%%%%%%%%%%%%%%%%%%%%%%%%%%%%%%%%%%%%%%%%

%%%%%%%%%%%%%%%%%%%%%%%%%%%%%%%%%%%%%%%
% Use English digits in equations
\usepackage{amsmath}

%\DefaultMathsDigits

% Default footnotes from left to right
\setfootnoteLR

% Use dash instead of dot in section numbers
%\SepMark{.}										


%% Change fonts and margins of section and subsection titles
%% For chapters please see firstpages.tex
%\titlespacing*{\section}{0pt}{1cm}{0.2cm}
%\titleformat{\section}
%{\fontsize{12}{6}\scshape\bfseries}{\thesection}{1em}{}
%
%\titlespacing*{\subsection}{0pt}{.8cm}{0cm}
%\titleformat{\subsection}
%{\fontsize{11}{6}\scshape\bfseries}{\thesubsection}{1em}{}
\begin{document}
% دستوری جهت حذف کردن شماره صفحه و سربرگ، در صورت وجود (فقط در صفحه جاری)
\thispagestyle{empty}
\begin{center}
	{\Naz
	باسمه تعالی
	\\[.1cm]
	\includegraphics[height=3cm]{iut_logo.png}
	%دستوری برای کم کردن فاصله بین لوگو و خط پایین آن
	\vspace{-2mm}
	
	{\large
	دانشگاه صنعتی اصفهان\\
	}
	{\small
	دانشکده برق و کامپیوتر
	}
	%دستوری برای تعیین فاصله بین دو خط
	\\[.3cm]
	{\large
		\textbf{پیشنهاد موضوع تحقیقاتی برای پایان‌نامه کارشناسی ارشد}	
		\\[.2cm]
	}

%%{
%%  \fontsize{9}{10.8}\selectfont دانشجوی محترم: لطفاً 
%%	\textbf{شیوه نامه تدوین گزارش پیشنهاد موضوع تحقیق پایان‌نامه کارشناسی ارشد}	را از وب سایت دانشکده دريافت نمایید.	نسخه 94/06/31
%%}

}

\end{center}
%%%%%%%%%%%%%%%%%%%%%%%%%%%%%%%%%%%%%%%%%%%
%\setlength{\parindent}{0pt}
%\vspace{-0.25cm}
%%%%%%%%%%%%%%مشخصات فردی را وارد کنید
{
\titr
\textbf{مشخصات دانشجو:}
}
{\Naz
\begin{center}
	\scalebox{1.1}{
		\begin{tabular}{|r|r|r|r|}
			
			\hline
			نام و نام خانوادگی: نوید نصیری&
			شماره دانشجويي: 40214394&
			رشته: مهندسی کامپیوتر&
			گرایش: نرم‌افزار
			\\ \hline
			\multicolumn{2}{|r|}{گرایش/رشته  دوره کارشناسی: مهندسی کامپیوتر-سیستم‌های هوشمند}&
			\multicolumn{2}{|r|}{شماره تلفن همراه: 09215312706}
			\\	\hline
			
		\end{tabular}}
	\end{center}}
%	\vfill

%%%%%%%%%%%%%%عنوان پروپوزال را وارد کنید
{
\titr
\textbf{مشخصات تحقیق پیشنهادی:}
}
{\Naz
\begin{center}
	\scalebox{0.97}{
    \begin{tabular}{|c|c|}
	\hline
	عنوان فارسی & دوقلو‌های مجازی چند عاملی در جهت کاهش هزینه‌ها
	\\ \hline  
	عنوان انگلیسی & \lr{Multi Agent Digital Twin To Reduce Costs}
	\\ \hline      
	\end{tabular}}
\end{center}
}
%%%%%%%%%%%%%% کلمات کلیدی گزارش
%کلمات کلیدی خود را وارد کنید.
\begin{center}
	\scalebox{1.}{
		\begin{tabular}{|r|r|r|r|}
			\hline
			کلمات کلیدی گزارش:&1- دوقلو‌های مجازی&2- سیستم‌ چند‌عاملی &3- هوش مصنوعی
			\\ \hline  
		\multicolumn{1}{|l|}{\lr{3- AI}} &
		\multicolumn{1}{|l|}{\lr{2- Multi Agent System}} &
		\multicolumn{1}{|l|}{\lr{1- Digital Twins}} &
		\multicolumn{1}{|l|}{Keywords:}
			\\ \hline
	\end{tabular}}
\end{center}


%%%%%%%%%%%%%%مشخصات استاد راهنما و مشاور را وارد کنید، موارد وارد شده برای مثال است.
{\titr
	\textbf{اساتید راهنما و مشاور:}}
\begin{center}
	\scalebox{0.89}{
		\begin{tabular}{|r|r|r|r|r|r|}
			\hline
			\multicolumn{1}{|c|}{سمت} &
			\multicolumn{1}{c|}{نام و نام خانوادگی} &
			\multicolumn{1}{c|}{درصد مشارکت} &
			\multicolumn{1}{c|}{دانشگاه/مؤسسه} &
			\multicolumn{1}{c|}{مرتبه علمی} &
			\multicolumn{1}{c|}{تخصص مرتبط با پایان نامه}
			\\ \hline
			\multicolumn{1}{|r|}{استاد راهنما} &
			دکتر  مسعودرضا هاشمی&۱۰۰ درصد&دانشگاه صنعتی اصفهان&استاد& شبکه‌های کامپیوتری
			\\ \hline
			\multicolumn{1}{|r|}{*استاد مشاور} &
			&&&&
			\\ \hline
			\multicolumn{1}{|r|}{*استاد راهنمای دوم} &
			&&&&
			\\ \hline
			\multicolumn{1}{|r|}{*استاد مشاور دوم} &
			&&&&
			\\ \hline
	\end{tabular}}
\end{center}
%	\vfill

\newpage
\setlength{\parindent}{0pt}
%\pagebreak
% دستوری جهت حذف کردن شماره صفحه و سربرگ، در صورت وجود (فقط در صفحه جاری)
\thispagestyle{empty}

\clearpage

	%%%%%%%%%%%سربرگ گزارش
	% دستوری جهت حذف کردن شماره صفحه و سربرگ، در صورت وجود (فقط در صفحه جاری)
\thispagestyle{empty}
\newgeometry{left=3cm,right=3cm,top=2cm}
{\Naz
\begin{center}
	باسمه تعالی
	\\[.1cm]
	\includegraphics[height=3cm]{iut_logo.png}
	%دستوری برای کم کردن فاصله بین لوگو و خط پایین آن
%	\vspace{-2mm}
	\vspace{0.2cm}
	
	{\large
		دانشگاه صنعتی اصفهان\\
	}
	{\small
		دانشکده برق و کامپیوتر
	}
	%دستوری برای تعیین فاصله بین دو خط
	\\[1.cm]
	{\LARGE
		\textbf{گزارش پيشنهاد موضوع تحقيق پایان‌نامه کارشناسی ارشد}	
		\\[1.3cm]
	}
%	در اینجا نسخه اول یا نهایی بودن گزارش خود را معین نمایید.
	{\Large
		\textbf{نسخه اول }
		\\[1.6cm]
	}
%	رشته خود را مشخص نمایید.
	{\LARGE
		\textbf{مهندسی کامپیوتر/نرم‌افزار}
		\\[1.2cm]
	}
%	عنوان را وارد کنید.
	{\titr
		{\LARGE
			\textbf{دوقلو‌های مجازی چند عاملی در جهت کاهش هزینه‌ها}
		\\[2.5cm]
		}
	}
%	نام خود را بنویسید
	{\LARGE
		\textbf{نوید نصیری}
		\\[2cm]
	}
%		نام استاد(های) راهنمای خود را وارد کنید. موارد نوشته شده برای مثال اند
	{\Large
		دکتر مسعودرضا هاشمی
%		\\
%		نام استاد مشاور را وارد کنید.
		\\[3.cm]
	}
	{\Large
	\textbf{تاریخ}
	\\[.4cm]
	\textbf{شهریور ۱۴۰۳}
	}
\end{center}}
\clearpage
% دستوری جهت حذف کردن شماره صفحه و سربرگ، در صورت وجود (فقط در صفحه جاری)
\thispagestyle{empty}
		

%*****************************************************************
%% تنظیم مناسب صفحه و فونت برای متن اصلی پایان‌نامه
\newgeometry{left=2.5cm,right=3cm,top=3cm,bottom=2.5cm,includehead=false,headsep=1cm,footnotesep=.5cm}
\settextfont{XB Zar}\fontsize{12}{6}\selectfont
\setlatintextfont{Times New Roman}\fontsize{11}{6}
\baselineskip=.9cm

\pagenumbering{arabic}
% Moving page number to top right
\pagestyle{myheadings}

\newpage

\section{موضوع کلی و زمینه اصلی تحقیق}

در دنیای امروز، با پیشرفت تکنولوژی، هزینه نگهداری و استفاده از آن نیز روز به روز افزایش می‌یابد. این افزایش هزینه مارا به بررسی روش‌ها و حالت‌های مختلف در دنیایی غیر از دنیای واقعی وا می‌دارد. ایجاد یک محیط مجازی از موجودیت‌های واقعی و فیزیکی یا به عبارتی دوقلوی مجازی\LTRfootnote{Digital Twin} 
می‌تواند این کاهش هزینه‌ را با خود به همراه داشته باشد. امروزه، مجازی‌سازی اشیاء، موجودیت‌ها و فرایند‌ها، این کمک را می‌کند که بتوان بدون نیاز به ایجاد تغییر و دستکاری در واقعیت، با مشاهده و بررسی دوقلوی مجازی آن، بهترین تصمیم گرفته و اجرا شود. با این وجود، یک حقیقت غیرقابل انکار این است که یک فرایند واقعی، از ارتباط و تاثیر موجودیت‌ها، اشیاء و فرایند‌های کوچک‌تر بر یکدیگر ایجاد می‌شود و نمی‌توان این مورد را نادیده گرفت. لذا مطالعه بر نوع ارتباط چند عامل\LTRfootnote{Multi Agent}
، مدیریت،‌ تبادل داده و تصمیم‌گیری درست بر اساس داده‌ها و هوش مصنوعی، می‌تواند به دوقلو‌های مجازی در دنیای مجازی کمک کند. پس دو مورد سیستم‌های چند‌عاملی و دوقلو‌های مجازی، می‌توانند مکمل یکدیگر باشند و نقاط ضعف یکدیگر را کامل کرده و در نهایت یک سامانه کامل از دنیای مجازی و دوقلوی مجازی به ما اهداء کنند.


\section{مرور کارها و نتایج  قبلی }

بر اساس 
\cite{b1}
دوقلو‌های مجازی معمولا از سه جزء ایجاد می‌شوند:\\
الف) فضای واقعی که دارای شئ واقعی است.\\
ب) فضای مجازی که دارای شئ مجازی است.\\
ج) ارتباط بین شئ مجازی و واقعی و انتقال داده میان آن‌ها به صورت بی‌درنگ که باعث می‌شود از سیستم حقیقی یک دوقلوی مجازی ایجاد شود.\\
در مقالات مختلف، ایجاد دوقلو‌های مجازی از اشیاء‌ها و موجودیت‌های مختلفی صورت می‌گیرد. طبق ادعایی که در مقاله 
\cite{b2}
صورت گرفته‌است، دوقلو‌های مجازی تنها یک شبیه‌سازی ساده نیست. بلکه یک مدل زنده، هوشمند و در حال تکامل است. دوقلوی مجازی، چرخه عمر همتای حقیقی خود را دنبال می‌کند تا بتوان آن را نظارت\LTRfootnote{Monitor}
و کنترل\LTRfootnote{Control}
کرد و همچنین کارکرد آن را بهبود بخشید. دوقلوی مجازی به صورت مستمر با فضای واقعی و محیط آن در ارتباط بوده و به صورت بی‌درنگ داده‌ بین یکدیگر ردوبدل می‌کنند.\\
یک نمونه از دوقلوی مجازی را می‌توان در مقاله 
\cite{b3}
مشاهده کرد که به یک شهر هوشمند\LTRfootnote{Smart City}
اشاره می‌کند. طبق ادعای مقاله، یک چارچوب\LTRfootnote{Framework}
برای ایجاد دوقلوی مجازی شهر هوشمند ایجاد شده‌است که با آن می‌توان به صورت کامل اطلاعات شهر را پایش کرده و شرایط آن را در زمان‌های مختلف مشاهده و پیش‌بینی کرد. این دوقلوی مجازی شهر هوشمند به کاربر این امکان را می‌دهد که با بررسی دقیق، برای شهر دقیق‌ترین تصمیم‌ها را گرفته\LTRfootnote{Decision Making}
 و اجرایی کند.\\
 مقاله 
 \cite{b4}
 به بررسی دوقلو‌های مجازی در اینترنت وسایل‌ نقلیه\LTRfootnote{Internet of Vehicles (IoF)}
 پرداخته است. در واقع این مقاله، یک سکو\LTRfootnote{Platform}
 جهت به اشتراک گذاری داده، بین وسایل نقلیه موجود در سطح شهر را معرفی می‌کند. در این مقاله، یک دوقلوی دیجیتال، به صورت یک بازتاب از وسیله نقلیه در فضای ابری تعریف می‌شود که داده‌ها را با وسیله نقلیه در هر لحظه همگام می‌کند. داشتن این دوقلوی مجازی در بستر ابری، باعث می‌شود که بتوان با به اشتراک‌گذاری داده‌ها بدون محدودیت فاصله و پهنای باند، کارایی رانندگی را افزایش داده و تصمیمات دقیق‌تری بر اساس داده‌ها گرفته‌ شود. این کاربرد نمونه‌ای از دوقلوی مجازی چند‌عاملی نیز به حساب می‌آید.\\
در نهایت مقاله 
\cite{b5}
یک چارچوب\LTRfootnote{Framework}
 به نام 
\lr{MADTwin}\LTRfootnote{Multi Agent Digital Twin}
ارائه می‌دهد که برای مدل‌سازی دوقلوی دیجیتال برپایه عامل‌\LTRfootnote{Agent}
مناسب است. از این چارچوب برای هوشمند‌سازی یک انبار‌ توسط ربات‌ها استفاده کرده تا بتوان به طور بهینه هر ربات را به سوی اهداف محلی خود و اهداف کل سیستم به صورت مشارکتی یا مستقل هدایت کند.



\section{نحوه ارزیابی دستاوردهای تحقیق}
بخش هوش مصنوعی جزء جدایی ناپذیری از دوقلو‌های مجازی است. یکی از روش‌های ارزیابی افزایش دقت خروجی‌های بر پایه روش‌های هوش مصنوعی، در مقایسه با کار‌های مشابه است. همچنین میزان اقتصادی بودن و قابل پیاده‌سازی بودن روش ارائه شده و در نهایت کاهش هزینه‌های کلی نیز هدف نهایی پروژه خواهد بود. 



\section{زمان‌بندی تحقیق}
در جدول \ref{tab1}، برنامه‌ی یکسال آینده، زمان‌بندی تقریبی انجام مراحل مختلف تحقیق جهت بدست آوردن نتایج تحقیق برآورد شده است.
\begin{table}[H]
	\centering
	\caption{پیش بینی زمان بندی فعالیت‌ها و مراحل اجرایی تحقیق و ارائه گزارش پیشرفت کار}
	\label{tab1}
	\scalebox{.80}{
		\begin{tabular}{|c|c|c|c|c|c|c|c|c|c|c|c|c|c|c|}
			\hline
			ردیف &
			\multicolumn{1}{>{\centering}m{7cm}|}{شرح فعالیت}
			&	 زمان کل &
			\multicolumn{12}{c|}{زمان اجرا به ماه}\\
			\cline{4-15}
			&
			\multicolumn{1}{>{\centering}m{4cm}|}{}&
			(ماه) & 1 & 2 & 3 & 4 & 5 & 6 & 7 & 8 & 9 & 10 & 11 & 12
			\\ \hline
			1 &
			\multicolumn{1}{>{\centering}m{7cm}|}{مطالعه و پیاده‌سازی کار‌های پیشین}	
			& 4 & $\times$ &$\times$ &$\times$ & $\times$& & & & & & & &
			\\ \hline
			2 &
			\multicolumn{1}{>{\centering}m{7cm}|}{بهبود روش‌های قبلی و ارائه راهکار جدید}	
			& 5 & & & $\times$ &$\times$ & $\times$&$\times$ &$\times$ & & & & &
			\\ \hline
			3
			&
			\multicolumn{1}{>{\centering}m{7cm}|}{شبیه‌سازی ساختار جدید و ارزیابی نتایج حاصله}
			& 5 & & &  & & & &$\times$ &$\times$ &$\times$ &$\times$ &$\times$ &
			\\ \hline
			4 &
			\multicolumn{1}{>{\centering}m{7cm}|}{نگارش پایان‌نامه}
			& 5 & & & &  &  &  &  & $\times$ & $\times$ & $\times$ & $\times$ & $\times$
			\\ \hline
%			5 &
%			\multicolumn{1}{>{\centering}m{7cm}|}{}
%			& & & & & & & & & & & & &
%			\\ \hline
	\end{tabular}}
\end{table}
\section{امكانات و تجهيزات مورد نياز، نحوه تامين و جدول  برآورد هزينه}


\newpage

\renewcommand{\bibname}{مراجع}
\addcontentsline{toc}{chapter}{مراجع}

\bibliographystyle{IEEEtran}

\setLTRbibitems

\begin{thebibliography}{99}

\begin{latin}

\bibitem{b1}
Grieves, M. and Vickers, J. (2017) 
\textit{‘Digital Twin: Mitigating Unpredictable, Undesirable Emergent Behavior in Complex Systems’}
, in Kahlen, J., Flumerfelt, S. and Alves, A. (eds) Transdisciplinary Perspectives on Complex Systems. Cham: Springer. Available at: \url{https://doi.org/10.1007/978-3-319-38756-7_4} 

\bibitem{b2}
Barricelli, B. R., Casiraghi, E. and Fogli, D. (2019) \textit{'A Survey on Digital Twin: Definitions, Characteristics, Applications, and Design Implications'}, IEEE Access, vol. 7, pp. 167653-167671. Available at: \url{https://doi.org/10.1109/ACCESS.2019.2953499}

\bibitem{b3}
Dani, A.A.H., Supangkat, S.H., Lubis, F.F., Nugraha, I.G.B.B., Kinanda, R. and Rizkia, I. (2023) \textit{‘Development of a Smart City Platform Based on Digital Twin Technology for Monitoring and Supporting Decision-Making’}, Sustainability, 15(18), 14002. Available at: \url{https://doi.org/10.3390/su151814002}

\bibitem{b4}
Tan, C., et al. (2023) \textit{‘Digital Twin Enabled Remote Data Sharing for Internet of Vehicles: System and Incentive Design’}, IEEE Transactions on Vehicular Technology, 72(10), pp. 13474-13489. Available at: \url{https://doi.org/10.1109/TVT.2023.3275591}

\bibitem{b5}
Marah, H. and Challenger, M. (2024) \textit{‘MADTwin: a framework for multi-agent digital twin development: smart warehouse case study’}, Annals of Mathematics and Artificial Intelligence, 92, pp. 975–1005. Available at: \url{https://doi.org/10.1007/s10472-023-09872-z}

\bibitem{b6}
Pretel, E., Navarro, E., López-Jaquero, V., Moya, A. and González, P. (2022) \textit{‘Multi-Agent Systems in Support of Digital Twins: A Survey’}, in Ferrández Vicente, J.M., Álvarez-Sánchez, J.R., de la Paz López, F. and Adeli, H. (eds) Bio-inspired Systems and Applications: from Robotics to Ambient Intelligence. IWINAC 2022. Lecture Notes in Computer Science, vol 13259. Springer, Cham, pp.  524–533. Available at: \url{https://doi.org/10.1007/978-3-031-06527-9_52}


\end{latin}

\end{thebibliography}


\newpage
	%%%%%%%% تعهدنامه


\begin{center}

	\scalebox{.80}{
		\begin{tabular}{|c|}
			\hline
			\\
			\multicolumn{1}{|>{\raggedright}m{19cm}|}{\textbf{تعهد دانشجو:}}
			\\
			\\
			%%%%%%%% اطلاعات (فرم را با اطلاعات خود پر کنید.)
			\multicolumn{1}{|>{\raggedright}m{19cm}|}{
				اینجانب ---------------- ضمن اطلاع از اینکه كليه حقوق مادی مترتب بر نتايج تحقيق این پایان نامه متعلق به دانشگاه صنعتی اصفهان بوده و انتشار نتايج تابع مقررات دانشگاهی است که صرفاً پس از موافقت و تایید استاد/ اساتيد راهنما قابل انجام است، تعهد می‌نمایم که اصول اخلاق علمی در انتخاب موضوع و نگارش این گزارش رعایت شده و در طول تحقیق و تدوین/انتشار  نتایج/گزارش ها/مقالات  بعدی نیز رعايت شود.
			}
			\\
			\multicolumn{1}{|>{\raggedright}m{19cm}|}{
				%				نام خود را وارد کنید.
				نام و نام خانوادگي: 
				\hspace{.5cm}
				امضا:
				%			فایل امضای خود را با نام Signature.JPG در مسیر قرار دهید و کامنت بعدی را انکامنت کنید
				%			\includegraphics[scale=.4]{Signature.JPG}
				\hspace{1.5cm}
				تاریخ:
			}
			\\ \hline
	\end{tabular}}
\end{center}
%\clearpage
%%%%%%%%%%%%%%%%%%%%%%%%%%%%%%%%%%%%%%%%%%%%%%%%%%%%%%%%
\begin{center}
	\scalebox{.80}{
		\begin{tabular}{|c|}
			\hline
			\\
			\multicolumn{1}{|>{\raggedright}m{19cm}|}{\textbf{تائيد اساتيد راهنما و مشاور پایان‌نامه:}}
			\\
			\\
		
			%%%%%%%% اطلاعات
			\multicolumn{1}{|>{\raggedright}m{19cm}|}{
				در صورت تصويب موضوع تحقيق پيشنهادي، بدين‌وسيله آمادگي خود را براي راهنمايي و مشاوره دانشجو در كليه مراحل انجام و ارائه تحقيق و مشاركت در ارزيابي پایان‌نامه بر اساس ضوابط دانشگاه و دانشكده اعلام مي‌نماييم. 
			}
			\\
			\multicolumn{1}{|>{\raggedright}m{19cm}|}{
				\hspace{5cm}نام و نام خانوادگی:
				\hspace{2.5cm}
				امضا:
				%			\begin{figure}
					%				\includegraphics[scale=0.4]{Signature.JPG}
					%			\end{figure}
				\hspace{2.5cm}
				تاریخ:
			}
			\\
			
			\multicolumn{1}{|>{\raggedright}m{18cm}|}{
				\begin{enumerate}[1-]
					\item
					استاد راهنما
					\item
					استاد مشاور 
					\item
					استاد راهنمای دوم*
					\item
					استاد مشاور دوم*
				\end{enumerate}
			}
			
			\\ \hline
	\end{tabular}}
\end{center}

%% نظر نهایی گروه

\begin{center}
	\scalebox{.80}{
		\begin{tabular}{|c|}
			\hline
			\\
			\multicolumn{1}{|>{\raggedright}m{19cm}|}{\textbf{نظر نهایی گروه:}}
			\\
			\\
			
			%%%%%%%% اطلاعات
			\multicolumn{1}{|>{\raggedright}m{19cm}|}{
				نسخه نهایی پیشنهاده مورد تایید است.
			}
			\\
			\multicolumn{1}{|>{\raggedright}m{19cm}|}{
				\hspace{5cm}نام و نام خانوادگی مدیر گروه:
				\hspace{2.5cm}
				امضا:
				%			\begin{figure}
					%				\includegraphics[scale=0.4]{Signature.JPG}
					%			\end{figure}
				\hspace{2.5cm}
				تاریخ:
			}
			\\
			
			\\ \hline
	\end{tabular}}
\end{center}

%% نظر نهایی شورای تحصیلات تکمیلی دانشکده

\begin{center}
	\scalebox{.80}{
		\begin{tabular}{|c|}
			\hline
			\\
			\multicolumn{1}{|>{\raggedright}m{19cm}|}{\textbf{نظر نهایی شورای تحصیلات تکمیلی دانشکده:}}
			\\
			\\
			
			%%%%%%%% اطلاعات
			\multicolumn{1}{|>{\raggedright}m{19cm}|}{
موضوع پیشنهادی تحقیق پایان‌نامه کارشناسی ارشد آقا/خانم --------------- مورد تایید است.
			}
			\\
			\multicolumn{1}{|>{\raggedright}m{19cm}|}{
				\hspace{5cm}نام و نام خانوادگی مدیر گروه:
				\hspace{2.5cm}
				امضا:
				%			\begin{figure}
					%				\includegraphics[scale=0.4]{Signature.JPG}
					%			\end{figure}
				\hspace{2.5cm}
				تاریخ:
			}
			\\
			
			\\ \hline
	\end{tabular}}
\end{center}


%\bibliography{REF}
\end{document}
%% End of file 